\documentclass[11pt]{article}
\usepackage[margin=1in]{geometry}
\usepackage{amsmath,amssymb,amsthm}
\usepackage[hidelinks]{hyperref}
\usepackage{enumitem}

\title{Literature Status Note: the Sub-Markovian Littlewood--Paley--Stein Question}
\author{Run \texttt{fa\_banach\_001}, lane 7}
\date{June 18, 2026}

\newcommand{\Pt}{P}
\newcommand{\Tt}{T}

\begin{document}
\sloppy
\maketitle

\paragraph{Status.}
\textbf{Literature already answered.}
The source question is from Teresa Mart\'inez, Jos\'e L. Torrea, and Quanhua Xu,
\emph{Vector-valued Littlewood-Paley-Stein theory for semigroups},
arXiv:math/0505303, Adv. Math. 203 (2006), 430--475.
The answering paper is Quanhua Xu,
\emph{Holomorphic functional calculus and vector-valued Littlewood-Paley-Stein
theory for semigroups}, arXiv:2105.12175.

\paragraph{The source question.}
In Section 3 of arXiv:math/0505303, Theorems 3.1 and 3.2
(\texttt{lusin semigroups} and \texttt{lusin type semigroups}) characterize
martingale cotype and martingale type by vector-valued
Littlewood--Paley--Stein inequalities for the Poisson semigroup
$\{\Pt_t\}_{t>0}$ subordinated to every symmetric diffusion semigroup
$\{\Tt_t\}_{t>0}$.
Immediately after the duality remarks following these theorems, the authors
write that all semigroups considered in that paper are Markovian
($\Tt_t1=1$), and that they do not know whether Theorems 3.1 and 3.2 still
hold for symmetric sub-Markovian semigroups, meaning the same assumptions
except Markovianity.

\paragraph{The answering theorem.}
Xu's 2021 paper explicitly identifies this as an open problem from the
earlier Littlewood--Paley--Stein papers. In the introduction it says that
Rota's dilation prevented weakening the symmetric diffusion hypothesis, that
Cowling's approach applies to symmetric submarkovian semigroups, and that it
had been open to establish the results of \cite{MTX2006} in Cowling's setting.
The paper then states that its objective is to resolve those problems.

The decisive result is Theorem 1.1, labelled \texttt{Poisson ML} in
arXiv:2105.12175.

\begin{quote}
Let $X$ be a Banach space and $1<p,q<\infty$. Let
$\{\Tt_t\}_{t>0}$ be a strongly continuous semigroup of regular operators on
$L_p(\Omega)$ and let $\{\Pt_t\}_{t>0}$ be its subordinated Poisson semigroup.
If $X$ has martingale cotype $q$, then $X$ has Luzin cotype $q$ relative to
$\{\Pt_t\}_{t>0}$. If $X$ has martingale type $q$, then $X$ has Luzin type
$q$ relative to $\{\Pt_t\}_{t>0}$.
\end{quote}

The same introduction notes immediately after the theorem that this improves
Theorem A from the earlier work because the semigroup is only assumed on a
single $L_p(\Omega)$, and markovianity, submarkovianity, and symmetry are not
assumed.

\paragraph{Why this answers the source question.}
A symmetric sub-Markovian semigroup is, in particular, a positive contraction
semigroup on the relevant $L_p$ spaces, and hence a semigroup of regular
operators on each fixed $L_p(\Omega)$. Applying Xu's Theorem \texttt{Poisson
ML} to such a semigroup gives the cotype and type inequalities for its
subordinated Poisson semigroup. Thus the ``if'' directions requested in the
sub-Markovian extension of Theorems 3.1 and 3.2 hold in a strictly broader
setting.

For the converse directions in the original ``for every semigroup'' formulation,
no new argument is needed: Markovian symmetric diffusion semigroups are a
subclass of symmetric sub-Markovian semigroups, so the original 2006 theorem
already forces martingale cotype/type when the inequality is required for every
symmetric sub-Markovian semigroup.

\paragraph{Scope.}
This note records the literature answer to the Poisson-subordinate
sub-Markovian extension of Theorems 3.1 and 3.2 from arXiv:math/0505303. It is
not a new proof or counterexample from this run. It does not settle the source
paper's separate Problem 1 on weak type $(1,1)$, nor all non-subordinated
heat-semigroup variants outside the hypotheses of Xu's later theorems.

\begin{thebibliography}{9}
\bibitem{MTX2006}
T. Mart\'inez, J. L. Torrea, and Q. Xu,
\emph{Vector-valued Littlewood-Paley-Stein theory for semigroups},
Adv. Math. 203 (2006), 430--475; arXiv:math/0505303.

\bibitem{Xu2021}
Q. Xu,
\emph{Holomorphic functional calculus and vector-valued Littlewood-Paley-Stein theory for semigroups},
arXiv:2105.12175.
\end{thebibliography}

\end{document}
